% ===================================================
% MỞ ĐẦU: THÔNG TIN NHÓM, PHÂN CÔNG VÀ THUẬT NGỮ
% Nguồn: GLOSSARY.md, docs/architecture/DEVELOPMENT_GUIDE.md,
%        docs/architecture/API_CONTRACTS.md, app/src/app.js
% ===================================================
\section*{Mở đầu: Giới thiệu, Phân công và Thuật ngữ Hệ thống}
\addcontentsline{toc}{section}{Mở đầu: Giới thiệu, Phân công và Thuật ngữ Hệ thống}

Báo cáo này trình bày quá trình phân tích yêu cầu, thiết kế hệ thống, hiện thực sản phẩm khả dụng tối thiểu (MVP) và kiểm thử cho đề tài \textbf{Smart E-Mobility Hub -- Hệ thống Điều phối Phương tiện Điện ĐHQG-HCM} trong học phần \textit{Công nghệ Phần mềm (CO3001 -- HK261)}. Nội dung báo cáo được tổ chức thành sáu chương bám sát các cột mốc của học phần: đặc tả yêu cầu và Use Case (Chương~1), mô hình động và thiết kế giao diện (Chương~2), kiến trúc 4+1 View, mô hình lớp và cơ sở dữ liệu (Chương~3), hiện thực MVP, kiểm thử và trình diễn (Chương~4), đánh giá kết quả và bài học kinh nghiệm (Chương~5), cùng khai báo sử dụng Generative AI và quản trị nhóm (Chương~6).

Về phạm vi kỹ thuật, phiên bản MVP sử dụng kiến trúc \textit{Modular Monolith} trên nền Node.js, Express và SQLite với ba phân hệ nghiệp vụ độc lập (\textit{Operations}, \textit{Mobility \& Simulation}, \textit{Parking \& Charging}). Các thao tác tương tác phần cứng hiện được mô phỏng thông qua bộ điều hợp giả lập trong tiến trình máy chủ; các giới hạn về phạm vi tính năng và những phần mới hoàn thiện ở mức dịch vụ được ghi nhận minh bạch tại các chương thiết kế, hiện thực và đánh giá.

\subsection*{Phân công Trách nhiệm Thành viên}

Nhóm triển khai song song theo bốn phân hệ độc lập, thống nhất giao tiếp qua hợp đồng REST API và giao diện nội bộ giữa các module. \Cref{tab:team-ownership} tóm tắt phạm vi phụ trách của từng thành viên.

\begin{table}[H]
\centering\small
\caption{Phân công trách nhiệm và tỷ lệ đóng góp của các thành viên}
\label{tab:team-ownership}
\begin{tabularx}{\textwidth}{>{\raggedright\arraybackslash}p{2.5cm} >{\raggedright\arraybackslash}p{2.6cm} >{\raggedright\arraybackslash}X >{\centering\arraybackslash}p{1.6cm}}
\toprule
\textbf{Phân hệ} & \textbf{Phụ trách (MSSV)} & \textbf{Phạm vi công việc chính} & \textbf{Đóng góp} \\
\midrule
Thiết kế UX/UI & Trần Đình Ý (2414100) & Thiết kế bộ 26 màn hình bản mẫu cho User Mobile, Operator Mobile và Operator Web Dashboard cùng luồng chuyển màn hình. & \textit{[Chờ XN]} \\
\addlinespace[4pt]
Thuê xe \& Mô phỏng & Nguyễn Đình Minh (2412079) & Hiện thực luồng tìm Hub, thuê/trả xe (\texttt{UC01--UC05}), mô phỏng \textit{What-if} (\texttt{UC13--UC14}), giao diện Live Dashboard và ứng dụng khách Android/Flutter. & \textit{[Chờ XN]} \\
\addlinespace[4pt]
Đỗ xe \& Sạc điện & Trần Văn Minh Triết (2413603) & Hiện thực luồng đỗ xe cá nhân (\texttt{UC06}), yêu cầu/dừng sạc (\texttt{UC07}, \texttt{UC07A}), lập lịch sạc (\texttt{UC12}) và lược đồ dữ liệu đỗ--sạc. & \textit{[Chờ XN]} \\
\addlinespace[4pt]
Vận hành Hub & Nguyễn Thanh Duy (2410536) & Hiện thực giám sát mạng Hub, quản trị tài nguyên, điều chuyển xe và xử lý sự cố (\texttt{UC08--UC11A}), cùng tầng truy cập SQLite dùng chung. & \textit{[Chờ XN]} \\
\bottomrule
\end{tabularx}
\end{table}

\subsection*{Bảng Thuật ngữ và Ký hiệu Chuyên ngành}

\Cref{tab:glossary} chuẩn hóa các thuật ngữ miền nghiệp vụ và kỹ thuật được sử dụng xuyên suốt báo cáo.

\begin{table}[H]
\centering\small
\caption{Thuật ngữ chuyên ngành và khái niệm miền nghiệp vụ cốt lõi}
\label{tab:glossary}
\begin{tabularx}{\textwidth}{>{\raggedright\arraybackslash}p{4.2cm} >{\raggedright\arraybackslash}X}
\toprule
\textbf{Thuật ngữ / Ký hiệu} & \textbf{Định nghĩa chuẩn hóa trong hệ thống} \\
\midrule
\textbf{User / Student} & Người dùng cuối (điển hình là sinh viên) tìm kiếm Hub, đặt, nhận/trả xe dùng chung, đặt chỗ đỗ và gửi yêu cầu sạc cho xe cá nhân. \\
\addlinespace[3pt]
\textbf{Operator} & Nhân viên vận hành giám sát mạng lưới Hub, quản trị tài nguyên, điều chuyển xe, xử lý sự cố, điều phối lịch sạc và chạy mô phỏng. \\
\addlinespace[3pt]
\textbf{Mobility Hub} & Trạm trung chuyển tại ĐHQG-HCM (Ga Metro ĐHQG, Ký túc xá khu A, ĐH Bách khoa -- Cơ sở 2, Thư viện Trung tâm), tích hợp xe dùng chung, chỗ đỗ và trụ sạc. \\
\addlinespace[3pt]
\textbf{Xe dùng chung / Xe cá nhân} & Phân loại quyền sở hữu phương tiện điện: xe dùng chung (\textit{Shared}) phục vụ cho thuê và điều chuyển; xe cá nhân (\textit{Private}) chỉ sử dụng dịch vụ đỗ và sạc. \\
\addlinespace[3pt]
\textbf{Cam kết đặt chỗ và Chiếm dụng vật lý} & Phân biệt giữa \textit{cam kết đặt chỗ} (phiếu \texttt{Reserved} theo hạn giữ hoặc khung giờ hẹn) với \textit{chiếm dụng vật lý} (bản ghi thực tế chỉ phát sinh sau khi check-in tại bãi). \\
\addlinespace[3pt]
\textbf{Phân bổ sạc và Phiên sạc} & Phân biệt giữa \textit{phân bổ cổng và công suất} (trạng thái \texttt{Reserved}) với \textit{phiên sạc thực tế} (trạng thái \texttt{Active}) chỉ bắt đầu sau khi thiết bị xác nhận cấp điện. \\
\addlinespace[3pt]
\textbf{Single-Winner Semantics} & Cơ chế bảo đảm chỉ một giao dịch thành công khi nhiều yêu cầu cùng tranh chấp một xe hoặc một chỗ đỗ thông qua khóa ghi nguyên tử trong cơ sở dữ liệu. \\
\addlinespace[3pt]
\textbf{Idempotency-Key} & Khóa định danh chống lặp gắn kèm các yêu cầu ghi HTTP, giúp phát lại kết quả an toàn khi mạng chập chờn mà không tạo giao dịch trùng lặp. \\
\addlinespace[3pt]
\textbf{Immutable Snapshot} & Bản chụp trạng thái toàn mạng lưới tại một thời điểm dùng riêng cho bộ máy mô phỏng \textit{What-if}, bảo đảm không làm biến đổi dữ liệu vận hành thực tế. \\
\bottomrule
\end{tabularx}
\end{table}
